\section{Protocol Record}
\label{app:protocol}

\paragraph{Registration} The ExCyTIn rules, the four hypotheses, the statistics, and the decision rule were written in a protocol file and committed together with the scripts before the single test run on the latest release (commit \texttt{0a22c60}; run committed as \texttt{4cd0f3a}). The hypotheses were: (H1) the graph lookup is exact on at least 25\% of test questions; (H2) at least half of the test questions name their end alert; (H3, exploratory) per model, the per-incident reported reward correlates positively with the per-incident named rate; (H4, exploratory) the alert-table lookup is within 10 points of the graph lookup. H1, H2, and H4 held; H3 was positive for 17 of 19 models on the latest release and for \oHPos{} of 19 on the reported release. The addendum for the reported release and the per-question logs (\cref{sec:protocol}) was committed (\texttt{50d6552}) before the logs were analyzed and run once (\texttt{a4c9672}). It registered: (A1) H1 and H2 on the reported release, which held; (A2, primary) at least 12 of 14 runs succeed more often on named questions, which did not hold (\lgPos{} of \lgModels); (A3) the same split by the alert-table lookup (\lgTPos{} of \lgModels); (A4, A5) the shortcut-resistant ranking and the share of successes the lookup also answers, both descriptive. The decision rule for A2 required us to drop any claim that agents' scores come from the shortcut, and the paper follows it. Commits are in the project's repository; they are not attested by an external registry.

\paragraph{Disclosure} Before this study, two feasibility scripts written for another purpose read all questions of the latest ExCyTIn release, including its 599 test questions, to check whether answers can be recovered from the graphs. The shortcut rules were written on the 418 training questions; the test questions were not used to write or tune them, but they had been read. The reported release and the agent logs were not read before the addendum, except one question and the field structure of one log file, which were needed to write it.

\paragraph{Exploratory analyses} The SIABench rule, the trace statistics of three runs (number of SQL actions, use of the alert table, wrong versus missing answers), and the split by path length were not registered. They are labeled as exploratory where they appear.

\section{ExCyTIn Details}
\label{app:excytin}

\paragraph{Baselines} Words are lowercase alphanumeric tokens longer than two characters, minus a stop list of 48 common and domain words (``alert'', ``related'', ``associated'', ``identify'', and similar). The requested entity type is the first match in this order: IP (``ip address''), hash (``sha256'', ``sha1'', ``hash''), URL (``url'', ``domain'', ``link''), mailbox (``email'', ``mailbox'', ``upn''), account (``account'', ``user'', ``sid''), file, process (``process'', ``command line''), host (``host'', ``device'', ``machine'', ``computer''). The alert-table lookup reads, for each type, the fields an analyst would report (for example \texttt{Address} for IPs; \texttt{Name}, \texttt{Sid}, \texttt{AadUserId}, or \texttt{UserPrincipalName} for accounts) and, among tied alerts, takes the earliest. Answers and gold values are lowercased and whitespace-normalized; a prediction is correct if it equals or contains the gold value.

\paragraph{Releases} The latest release (\texttt{opus5}) has 418 training and 599 test questions; the reported release (\texttt{o1}, test only for evaluation) has 589 test questions in two versions, the original used for the published scores and a corrected one. The lookups give \oGraph\% (graph) and \oTable\% (table) on the original and nearly the same on the corrected version.

\begin{table}[h]
  \caption{Reported release: per-incident question counts, named and lookup rates (\%), and the reported mean reward (\%) of three models~\citep{excytin}.}
  \label{tab:incidents}
  \centering\footnotesize
  \begin{adjustbox}{max width=\columnwidth}% generated by paper/make_audit.py -- do not edit
\begin{tabular}{@{}lrrrrrrr@{}}
\toprule
Incident & $n$ & Named & Graph & Table & GPT-4o & o3 & Claude-Opus-4.5 \\
\midrule
5 & 98 & 29.6 & 29.6 & 36.7 & 33.8 & 43.3 & 56.7 \\
34 & 82 & 63.4 & 31.7 & 31.7 & 29.3 & 35.9 & 75.4 \\
38 & 11 & 45.5 & 36.4 & 36.4 & 36.4 & 72.7 & 76.4 \\
39 & 98 & 54.1 & 24.5 & 25.5 & 27.3 & 44.3 & 49.1 \\
55 & 100 & 59.0 & 20.0 & 13.0 & 24.9 & 45.9 & 51.4 \\
134 & 57 & 73.7 & 57.9 & 22.8 & 49.1 & 63.2 & 82.5 \\
166 & 87 & 57.5 & 36.8 & 3.4 & 16.6 & 37.9 & 54.5 \\
322 & 56 & 83.9 & 35.7 & 35.7 & 31.5 & 54.3 & 66.4 \\
\bottomrule
\end{tabular}
\end{adjustbox}
\end{table}

\begin{table}[h]
  \caption{Reported release: per-model Spearman correlation over the eight incidents between the reported reward and the share of named questions, with the reported mean reward (\%).}
  \label{tab:h3}
  \centering\scriptsize
  \begin{adjustbox}{max width=\columnwidth}% generated by paper/make_audit.py -- do not edit
\begin{tabular}{@{}lrr@{\hspace{1em}}lrr@{}}
\toprule
Model & Reward & $\rho$ & Model & Reward & $\rho$ \\
\midrule
Claude-Opus-4.5 & 60.6 & +0.31 & GPT-4o & 29.3 & +0.02 \\
GPT-5.1 & 58.2 & +0.38 & Llama4-17b-Mav & 29.0 & +0.17 \\
Claude-Sonnet-4.5 & 48.7 & +0.21 & GPT-4.1-mini & 27.1 & +0.46 \\
o3 & 45.6 & +0.12 & Llama4-17b-Scout & 26.2 & +0.81 \\
o4-mini & 36.8 & +0.36 & o1-mini & 22.2 & +0.88 \\
Grok-4 & 34.4 & +0.29 & GPT-4o-mini & 19.2 & +0.54 \\
GPT-4.1 & 33.8 & +0.33 & Qwen-3-32b & 18.2 & +0.60 \\
Gemini 2.5 Flash & 30.5 & +0.45 & GPT-4.1-nano & 13.6 & +0.19 \\
Qwen-3-235b-thinking & 30.2 & +0.19 & Phi-4-14B & 8.5 & $-$0.07 \\
o3-mini & 29.6 & +0.38 & & & \\
\bottomrule
\end{tabular}
\end{adjustbox}
\end{table}

\section{GUIDE Details}
\label{app:guide}

We aggregate the 9.5 million evidence rows of the training split into 707{,}108 graded incidents (no incident has conflicting grades) and never read the columns that record the response taken after triage. The history rule uses, for each organization, its earliest 80\% of incidents in time to compute the most frequent grade per detector and predicts the later 20\% (\gDev{} incidents). The organization-disjoint comparison splits organizations 80/20 and scores naive Bayes classifiers over ten metadata facets (alert category, ATT\&CK technique, entity types, roles, scale, geography, automated-investigation verdict, detector history across organizations, and others) on 10{,}000 incidents of the held-out organizations, querying 1 to 12 cost units of facets in information-gain, fixed, or random order; all reach \gHeldMin--\gHeldMax{} accuracy against a majority class of \gHeldMajority.

\section{OTRF Details and the Disguised Task}
\label{app:otrf}

We downloaded 27 recordings of remote-execution techniques (remote services, WMI, scheduled tasks, DCOM, PowerShell remoting). Alerts are service installations (events 4697, 7045) and scheduled-task creations or updates (4698, 4702), deduplicated by name; an alert is labeled an attack if its name appears as a word in the recording's own metadata. The disguised task, \texttt{EventCacheManager} in the folder \texttt{\textbackslash Microsoft\textbackslash Windows\textbackslash Software\-Protection\-Platform}, appears in two recordings that reference Microsoft's analysis of the Solorigate intrusion~\citep{solorigate}. Its folder and name match genuine Windows tasks, so the text rule calls it benign. The creating account separates it: the domain user \texttt{pgustavo} created it over a network logon from another host, whereas every other task under \texttt{\textbackslash Microsoft\textbackslash Windows\textbackslash} in the recordings was created by a machine account.

\section{Datasets and Licences}
\label{app:licences}

ExCyTIn-Bench: code MIT, data CDLA-Permissive-2.0, released for research; question sets \texttt{opus5} and \texttt{o1}, the investigation graphs, the \texttt{SecurityAlert} tables of the public data release, and the published agent logs in \texttt{latest\_experiments} of the \texttt{microsoft/SecRL} repository (accessed 2026-09-24 and 2026-10-01); we keep only the question, gold answer, and reward of each log entry, and, for three runs, action counts. SIABench: the alert-triage set (35 JSON files) from the authors' repository, built from CIC-IDS2017, whose terms permit research use; we did not download the packet captures. GUIDE: anonymized incidents released by Microsoft; we read only the training split and do not redistribute per-incident data. OTRF Security-Datasets: MIT licence (LICENSE file); lab recordings without personal data.
